
\section{Problem Statement and Motivation}

Consider a conductive element with length $L$, cross-sectional area $A$, and resistivity $\rho$. 
Its electrical resistance is given by

\begin{equation}
    R = \frac{\rho L}{A}.
\end{equation}

Under mechanical loading, both geometry and material properties change. In particular, 
the length, area, and resistivity become functions of deformation, making resistance 
dependent on the mechanical state.

Using Ohm's law,

\begin{equation}
    R = \frac{V}{I},
\end{equation}

and expressing voltage and current as

\begin{equation}
    V = EL, \qquad I = JA,
\end{equation}

we obtain

\begin{equation}
    \frac{EL}{JA} = \frac{\rho L}{A}.
\end{equation}

Canceling geometric terms yields the local relation

\begin{equation}
    \mathbf{J} = \sigma \mathbf{E}, \qquad \sigma = \frac{1}{\rho},
\end{equation}

where $\mathbf{E} = -\nabla V$.

Combining this with charge conservation in steady state,

\begin{equation}
    \nabla \cdot \mathbf{J} = 0,
\end{equation}

leads to the conductivity equation

\begin{equation}
    \nabla \cdot (\sigma \nabla V) = 0.
\end{equation}

The resulting equation is an elliptic PDE in divergence form, belonging to a well-studied class 
of second-order equations in the theory of partial differential equations \cite{evans}.

In piezoresistive materials, the conductivity depends on strain:

\begin{equation}
    \sigma = \sigma(\varepsilon(u)),
\end{equation}

where the strain tensor is given by\footnote{The strain tensor is defined as the symmetric part of the displacement gradient, which removes rigid body rotations and captures the true deformation of the material.}

\begin{equation}
    \varepsilon(u) = \frac{1}{2}(\nabla u + \nabla u^T).
\end{equation}

This introduces a nonlinear coupling between mechanical deformation and electrical conduction.


\begin{figure}[htbp]
    \centering
    \includegraphics[width=0.7\textwidth]{images/model1.png}
    \caption{Change of geometric and material properties under mechanical deformation. $L$ denotes the length of the conductor, $W$ its characteristic width, 
and $A = W^2$ the cross-sectional area (assuming a square cross-section). 
The material property $\rho$ represents the electrical resistivity, which may vary under mechanical deformation.
    Both the dimensions and the resistivity of the conductor vary, leading to a strain-dependent resistance.}
\end{figure}

This illustrates that both geometric deformation and intrinsic material changes contribute to the variation of resistance.